\begin{marginfigure}[-1cm]
    \includegraphics[width=5cm]{images/jacques_hadamard.jpg}
    \caption{Jacques \textsc{Hadamard}}
\end{marginfigure}

\begin{tcolorbox}
    Soit $M \in \M_n(\C)$ et soient $X_1, \dots, X_n$ les vecteurs colonnes de $M$. Alors,
    $$|\mathrm{det}(M)| \leqslant \prod_{i=1}^{n} \Vert X_i \Vert$$
    avec égalité si et seulement si la famille $(X_i)$ est orthogonale.
\end{tcolorbox}

\begin{marginfigure}[0.8cm]
    \resizebox{6.5cm}{6.5cm}{%
    \begin{tikzpicture}
        \node[block] (gram) {Déterminant \\ de \textsc{Gram}};
        \node[block, right of=s1] (iwasama) {Décomposition \\ d'\textsc{Iwasama}};
        \node[block, below right of=gram] (hadamard) {Inégalité \\ d'\textsc{Hadamard}};
    
        \draw (gram) edge[bend right, above left] node {...} (hadamard);
    
        \draw (iwasama) edge[bend left] node {$M=OT$} (hadamard);
    
    \end{tikzpicture}    
}
\end{marginfigure}


\begin{preuve_titre}[par le déterminant de \textsc{Gram}]
    d
\end{preuve_titre}

\begin{remarque}
    L'inégalité de \textsc{Hadamard} nous apprend en fait que le volume du parallélotope défini par les vecteurs colonnes est inférieur ou égal au produit des normes de ses vecteurs et il y a égalité si et seulement si la matrice est diagonale, ou encore que le parallélotope est rectangle. 
\end{remarque}

\begin{tcolorbox}
    Soient $\mathscr{S}_n ^{++} (\R)$ l'ensemble des matrices réelles d'ordre $n$ symétriques à valeurs propres strictement positives et $A = (a_{i,j}) \in \mathscr{S}_n ^{++} (\R)$. Alors,
    $$\det(A) \leqslant \prod_{i=1}^{n} a_{i,i}.$$
\end{tcolorbox}

\underline{Démonstration:} (exercice 4, TD 14 \cite{acamanes})

\begin{enumerate}
    \item Soit $(\gamma_1, \dots, \gamma_n) \in (\Re)^n$. Montrer que $B = (\gamma_i \gamma_j a_{i,j}) \in \mathscr{S}_n^{++}(\R)$. 
    \item Montrer que $\det(A)^{1/n} \leqslant \frac{\Tr(A)}{n}$. \\
    \emph{On pourra utiliser l'inégalité arithmético-géométrique}.
    
    \marginnote[-2cm]{
    	\begin{kaobox}[frametitle=Inégalité arithmético-géométrique]
            Soient $n \in \Ne$ et $x_1, \dots, x_n$ des réels positifs. Alors, 
            $$\frac{x_1 + \cdots + x_n}{n} \geqslant \sqrt[n]{x_1 \cdots x_n}.$$
            Il y a égalité si et seulement si tous les $x_i$ sont égaux.
        \end{kaobox}
        Pour d'autres inégalités, lire le chapitre 16, p.117 de la deuxième édition de Raisonnements divins (en fr)
    }
    
    \item Montrer que pour tout $i \in \llbracket 1, n \rrbracket, a_{i,i} > 0$. On pose $\gamma_i = \frac{1}{\sqrt{a_{i,i}}}$. En déduire l'inégalité d'\textsc{Hadamard}.
\end{enumerate}